\section{Connections to Prior Work}
\label{sec:related}

We have demonstrated in Section \ref{sec:hubert} how our framework
instantiates HuBERT.
Given the generality of our framework, we can also instantiate other objectives that are similar to other self-supervised objectives.
In this section, we will discuss what our framework can achieve and how it differs from prior work,
including approaches based on likelihood and contrastive learning.
We do not consider combined objectives in this section, such as
w2v-BERT \citep{chung2021w2v}.
%
%

\subsection{APC and VQ-APC}

We begin with other variants of predictive coding for speech representation that optimizes the likelihood
\eqref{eq:pc}, in which $x_A$ is $x_{<t}$, $x_B$ is $x_{\geq t}$.
First, autoregressive predictive coding \citep{chung2019unsupervised,yang2022autoregressive} (APC) is a generalization of linear predictive coding \citep{atal1971speech,saito1967theoretical} (LPC), 
where a model is trained to predict future frames given the past.
Its vector-quantized variant, VQ-APC \citep{chung2020vector}, 
includes a quantization layer in the neural network while optimizing APC.
Both APC and VQ-APC optimize the likelihood, while our framework is based on the variational bound in equation \ref{eq:vlb}.
APC and VQ-APC in their original form in \citet{chung2019unsupervised} and \citet{chung2020vector} do not have the latent variables explicitly stated.
Our framework makes the choice of latent variables explicit, and the proposed auxiliary distribution 
offers more flexibility to estimate the likelihood, especially when marginalization of 
latent variables is not tractable.

Similar to our approach, \citet{yang2022autoregressive} and \citet{yeh2022autoregressive} make the latent variables explcit, leading to various other interpretations of APC with respect to, e.g., mutual information and co-training \citep{mcallester2018information}.

\subsection{MPC and DeCoAR}

MPC \citep{jiang2019improving,zhang2021transformer} and DeCoAR \citep{9053176} are both generalizations of APC,
replacing future prediction with masked prediction.
DeCoAR 2.0 \citep{ling2020decoar} adds quantization similar to VQ-APC.
Both optimize the likelihood in equation \ref{eq:pc}.
The difference between them lies in how the frames are masked---MPC masks many small spans while DeCoAR masks a single large span.

\subsection{HuBERT, WavLM, and BEST-RQ}
\label{sec:hubert-variants}

We have shown that HuBERT is a special case of our framework.
%
%
Subsequent work, such as WavLM \citep{chen2022wavlm} and BEST-RQ \citep{chiu2022self}, uses the same HuBERT objective.
WavLM improves over HuBERT with data augmentation.
BEST-RQ avoids updating the codebook, but requires a careful initialization.

When training HuBERT in \citet{hsu2021hubert}, there are multiple iterations, each of which trains a Transformer encoder from scratch.
What differs for each iteration are the training targets.
In the first iteration, quantized MFCC is used as targets, while in the second iteration, quantized hidden vectors from layer 9 of the first iteration are used as targets.
Our framework can also instantiate the second iteration with a pre-trained Transformer encoder from the first iteration. 
We can simply let
\begin{align}
    q(z_i|h_i) = \frac{\exp\left(-\|h_i - v_{z_i}\|^2 / \tau\right)} 
  {\sum_{k=1}^{K} \exp\left(-\|h_i - v_k\|^2 / \tau\right)},
\label{eq:softmax-q-iter2}
\end{align}
where the $i$-th frame $h_i$ now comes from the intermediate layer (e.g., 6th layer) of
the encoder from the first iteration.

\subsection{CPC, VQ-CPC, and wav2vec 2.0}
\label{sec:wav2vec}

%
%
%
%
%
%
CPC \citep{oord2018representation}, its vector-quantized variants, VQ-CPC \citep{niekerk2020vector}, 
wav2vec \citep{schneider2019wav2vec}, and wav2vec 2.0 \citep{baevski2020wav2vec}, 
are all based on noise contrastive estimation (NCE) \citep{smith-eisner-2005-contrastive,pmlr-v9-gutmann10a}.
%
%
%
%
%
In this section, we show how wav2vec 2.0 can be instantiated with our framework.

We derive their contrastive objective in wav2vec 2.0 from the cross entropy in the variational objective
 $\mathbb{E}_{z_i \sim q}[-\log p(z_i | x_{\setminus M})]$.
Recall that wav2vec 2.0 has a CNN followed by VQ and a Transformer,
where the VQ produces targets for contrastive learning.
%
We first choose
\begin{align}
q(z_i | x_M) &= \delta_{z_i=g(x_M)} \label{eq:w2v2-q}\\
p(z_i | x_{\setminus M}) &= \frac{\exp(\mathrm{sim}(\mathrm{enc}(x_{\setminus M})_i, z_i))}{
  \int \exp(\mathrm{sim}(\mathrm{enc}(x_{\setminus M})_i, z)) dz},
\end{align}
where $g(\cdot)$ is the CNN with the VQ, $\mathrm{enc}(\cdot)$ is the Transformer,
$\delta$ being the Dirac delta function, and $\mathrm{sim}(x, y) = \cos(x, y)$.
The choice of $p(z_i | x_{\setminus M})$ is similar to that of \Eqref{eq:prior} in our framework,
except a similarity function and a continuous $z$.
The cross entropy becomes
\begin{align}
    \mathbb{E}_{z_i \sim q}&[-\log p(z_i | x_{\setminus M})] \notag\\
    =&\log\frac{\exp(\mathrm{sim}(\mathrm{enc}(x_{\setminus M})_i, z_i))}{
   \int\exp(\mathrm{sim}(\mathrm{enc}(x_{\setminus M})_i, z')) dz}.
\end{align}
Note that $z_i$ here is the codeword after quantization, i.e., a continuous vector \citep{baevski2020wav2vec}.
The denominator in the cross entropy is in general difficult
to compute due to the integral.\footnote{In fact,
even when $\mathrm{sim}(x, y) = x^\top y$, i.e.,
$p(z_i | x_{\setminus M})$ being a von-Mises--Fisher distribution,
the denominator of $p$ is still difficult to compute.}
Inspired by NCE, \citet{oord2018representation} and \citet{baevski2020wav2vec}
approximate the integral by summing over a set of negative samples $S_i$. 
The cross entropy becomes
\begin{align}
    \log \frac{\exp(\mathrm{sim}(\mathrm{enc}(x_{\setminus M})_i, z_i))}{
  \sum_{z' \in S_i} \exp(\mathrm{sim}(\mathrm{enc}(x_{\setminus M})_i, z'))}.
  \label{eq:w2v2}
\end{align}
which is the objective used in wav2vec 2.0.
Consequently, minimizing the contrastive loss is analogue to minimizing the cross entropy between 
the masked frames and unmasked frames after quantization.

%
%
%
%
%
%
%
%
%
%
%
%
%
%
%
%
%
%
%
%
%
%
%

%
%
%
%
%
%
%
Several design choices in wav2vec 2.0 are taken verbatim by HuBERT, but they have a few key differences.
Both use a softmax to parameterize $p(z_i | x_{\setminus M})$, but wav2vec 2.0 assumes $z_i$ to be a continuous vector, while HuBERT assumes $z_i$ to be discrete.
As a consequence, wav2vec 2.0 uses a softmax over the negative and positive samples, while HuBERT uses a softmax over the codewords in a codebook.
The codebook in wav2vec 2.0 is trained together with the rest of the model.
In other words, there is no reconstruction and we can simply parameterize
$p(x_i | z_i)$ as a uniform distribution. 
In contrast, the codebook in HuBERT is trained with an offline $k$-means.
Compared to our framework, we have the choice to optimize the codebook
offline or jointly with the Transformer.

There is a additional diversity loss in wav2vec 2.0 that
computes the entropy of codeword usage within a batch,
to encourage the use of individual codewords.
The diversity loss was not meant to be part of the main objective
\cite{baevski2020wav2vec}, but more of a regularizer that improves
training.
However, our instantiation on HuBERT naturally has an entropy term
per frame that serves a similar purpose as the diversity loss.

%
%

%
%
%
%
%
%
%
%
%
%
%
%
%
%
%
%
%
%
%
